\section{Physical Background}

\subsection{Equations of Motion}

The spinning motion of the rigid system can be described using the Euler angles $\theta$, $\psi$, and $\phi$:
\begin{itemize}
    \item $\theta:=$ Tilt angle with the horizontal plane.
    \item $\psi:=$ Precession angle of the contact point $P$.
    \item $\phi:=$ Spin angle in the disk symmetry axis.
\end{itemize}
\begin{equation}
    \vec\omega = \dot\psi\hat{z}+\dot\theta\hat{e}_{2}+\dot\phi\hat{e}_{3}
\end{equation}
where $\omega$ is the angular velocity vector of the disk, $\dot\psi$ is the precession rate, $\hat{e}_{2}$ is the vector orthogonal to $\hat{z}$ and $\hat{e}_{3}$, and $\hat{e}_{3}$ is the symmetry axis of the disk. The nutation angle $\theta$ has the value $\theta=0$ for a flat disk and $\theta=\pi/2$ when it is upright. 

The non-slip condition in the precession contact point $P$ is $\dot r_{P}=0$, thus we can decompose it into:
\begin{equation}
    \dot{\vec{r}}_{P}=\dot{\vec{r}}_{cm}+\vec{\omega}\times\vec{r}_{cm\rightarrow P}=0
\end{equation}

We will define the precession rate as $\Omega:=\dot\psi$.
Given $I_{\parallel}=\frac{1}{2}mR^2$ and $I_{\perp}=\frac{1}{4}mR^2$, the angular momentum $\vec L=I_{\parallel}\omega\hat e_r + I_{\perp}\Omega\sin\theta\hat e_z + ...$ can be expressed as

\begin{equation}
    \vec L=\frac{1}{2}mR^2\omega\hat e_r + \frac{1}{4}mR^2\Omega\sin\theta\hat e_z
\end{equation}
\begin{equation}
    \dot{\vec L}=\tau=mgR\sin\theta\hat e_{\tau}
\end{equation}
We can arrive then to the key relation between the tilt angle and $\Omega$:
\begin{equation}
    \quad \Omega^2=\frac{4g}{R\sin\theta}
\end{equation}

Having defined the energy of the system as $E=\frac{3}{2}mgR\sin\theta$ and considering the small angle approximation, the potential $\Phi$ can be derived as
\begin{equation}
    E=-\nabla \Phi \quad \rightarrow \quad \Phi=-\frac{3}{2}mgR\dot\theta  
\end{equation}
\begin{equation}
    \dot\theta=-\frac{2\Phi}{3mgR}
    \label{eq:tilt-derivative}
\end{equation}

\subsection{Dissipation Potentials}
The two main dissipation components of the system are the rolling friction between the disk and the surface material where it spins and the drag force experimented in the air.
\subsubsection{Rolling friction regime}
The potential of the system in the rolling friction regime is:
\begin{equation}
    \Phi_{roll}=\eta m g R \cos{\theta}\Omega\approx\eta mgR\Omega
\end{equation}
Substituting in Equation~\ref{eq:tilt-derivative} results in the following differential equation:
\begin{equation}
    \dot\theta\approx-\frac{2\eta}{3}\Omega=-\frac{2\eta}{3} \left( \frac{4g}{R} \right)^{1/2}\theta^{1/2}
\end{equation}
Upon solving this ODE, we can arrive to the power term for the rolling friction dissipation $n_{\theta}=2/3$:
\begin{equation}
    \theta\propto(t_f-t)^{2/3}
\end{equation}

\subsubsection{Air friction regime}
The drag potential of the disk inside a fluid such as air is determined on the fluid model that one uses. A simpler model\cite{Moffatt2000} can be developed considering air as a lubrication layer of thickness $h\sim R\theta$, thus
\begin{equation}
    \Phi_{moffet}=\pi\eta gR^2\theta^{-2}
\end{equation}
\begin{equation}
    \dot\theta_{moffet}=\frac{-2\pi\eta R}{3m\theta^2}
\end{equation}
Upon solving this ODE, we can arrive to the power term for the rolling friction dissipation $n_{\theta}=1/3$:
\begin{equation}
    \theta_{moffet}\propto(t_f-t)^{1/3}
\end{equation}

However, if we instead consider a viscous layer of $\delta=\sqrt{\frac{2\eta}{\rho\Omega}}$, we can define the potential as:
\begin{equation}
    \Phi_{BL}=C\eta\left( \frac{R\Omega}{\delta} \right) R^2\delta
\end{equation}
\begin{equation}
    \Phi_{BL}=5g^{5/4}R^{11/4}\sqrt{\eta\rho}\theta^{-5/4}
\end{equation}
Which leads to the following ODE
\begin{equation}
    \dot\theta_{BL}=-\frac{8g^{1/4}R^{7/4}\sqrt{\eta\rho}\theta^{-5/4}}{3m}
\end{equation}
Upon solving this ODE, we can arrive to the power term for the rolling friction dissipation $n_{\theta}=4/9$:
\begin{equation}
    \theta_{BL}\propto(t_f-t)^{4/9}
\end{equation}

\subsection{Observables}
The Euler disk is a rigid body which rotates on a flat surface under the influence of gravity and dissipative forces. The rotation is started by the experimenter and is affected by the starting conditions. The motion of the disk can be analyzed as a combination of rotation and precession. As the disk loses energy due to dissipation forces (such as the air viscosity and rolling/sliding frictions) the inclination angle decreases while the precession frequency increases.

One of the best-known phenomena of the Euler disk system is that the system acts as finite-time singularity as the precession frequency $\Omega(t)$ follows a power-law behavior of the form \cite{Darmendrail2021}:

\[
\Omega^{-2} \propto \theta=A(t_f - t)^{-n}
\]

Where $t_f$ denotes the time the disk finally stops, $t$ is the actual time and $n$ is an exponent which is determined by the dissipation mechanisms.

Furthermore, the dominant acoustic frequency measured $f(t)$ is expected to be proportional to the precession frequency $\Omega(t)$, as suggested by the previous experimental studies using acoustic measurements.

This relation can be expressed as
\[
\Omega(t) \propto 2\pi f(t) 
\]


In this project, we investigate the behavior of the system explained above experimentally using acoustic measurement and spectral analysis with help of the Fourier Transformation.


\section{Measured Quantities}

The main measured quantity is the acoustic signal emitted by the disk which is caught by the phone's microphone. From the said signal, the dominant frequency $f(t)$ will be extracted with the help of Fourier analysis.

The following quantities will be determined:

\begin{itemize}
    \item Dominant frequency $f(t)$
    \item Stopping time $t_f$
    \item Dynamical exponent $n$ from fitting the measured frequency evolution of $f(t)$ to a power law of the form:
    \[
    f(t) \propto (t_f - t)^{-n}.
    \]
\end{itemize}

